%!TeX program = xelatex
% 使用 SYSUReport 模板，作者 NorthSecond，CC BY 4.0，保留模板原始署名。
\documentclass{SYSUReport}
\usepackage{array,tabularx,longtable,placeins,float,xurl,caption,etoolbox,needspace}
\captionsetup[figure]{labelformat=empty}
\usepackage{tikz}
\usetikzlibrary{arrows.meta,positioning,fit}
\setmainfont{Times New Roman}
\setCJKmainfont[AutoFakeBold=2,AutoFakeSlant=0.15]{SimSun}
\setCJKsansfont{SimHei}
\setCJKmonofont{SimSun}
\setmonofont{Consolas}
\setlength{\headheight}{16pt}
\setlength{\parskip}{0.22em}
\renewcommand{\arraystretch}{1.25}
\setcounter{tocdepth}{2}
\BeforeBeginEnvironment{lstlisting}{\par\vspace{0.8\baselineskip}\noindent\begin{minipage}[t]{\linewidth}\vspace{0.3\baselineskip}}
\AfterEndEnvironment{lstlisting}{\end{minipage}\par\vspace{0.6\baselineskip}}
\newcolumntype{P}[1]{>{\raggedright\arraybackslash}p{#1}}
\newcolumntype{Y}{>{\raggedright\arraybackslash}X}
\newcommand{\code}[1]{\begingroup\Urlmuskip=0mu plus 1mu\expandafter\nolinkurl\expandafter{\detokenize{#1}}\endgroup}
\newcommand{\pathCN}[1]{\texttt{\detokenize{#1}}}
\newcommand{\shot}[4][0.48\textheight]{\begin{figure}[H]\centering\includegraphics[width=\linewidth,height=#1,keepaspectratio]{#2}\caption{#3}\label{#4}\end{figure}}
\lstset{basicstyle=\ttfamily\small,keywordstyle=\color{myteal},commentstyle=\color{gray},stringstyle=\color{black},breaklines=true,columns=fullflexible,keepspaces=true,showstringspaces=false,frame=single,rulecolor=\color{gray!35},backgroundcolor=\color{gray!4},xleftmargin=0.3em,xrightmargin=0.3em,aboveskip=0.5em,belowskip=0.5em}
\headl{数据库原理实验}
\headc{}
\headr{第5--9周阶段报告}
\fancypagestyle{plain}{\fancyhf{}\fancyhead[L]{\fangsong 数据库原理实验}\fancyhead[R]{\fangsong 第5--9周阶段报告}\fancyfoot[C]{\thepage}}
\lessonTitle{数据库原理实验报告}
\reportTitle{二次元周边商品店数据库重构与应用实现}
\stuname{郑立峰\quad 李维政}
\stuid{24325341\quad 24325143}
\inst{计算机学院}
\major{计算机科学与技术}
\date{2026年10月24日}
\hypersetup{hypertexnames=false,pdftitle={二次元周边商品店数据库重构与应用实现 第5至9周阶段报告},pdfauthor={郑立峰 李维政}}
\begin{document}
\hypersetup{pageanchor=false}
\cover\thispagestyle{empty}\clearpage
\pagenumbering{Roman}\setcounter{page}{1}\hypersetup{pageanchor=true}
\begingroup\fontsize{9.5pt}{10.8pt}\selectfont\setlength{\parskip}{0pt}\tableofcontents\endgroup
\clearpage\pagenumbering{arabic}\setcounter{page}{1}

\section{实验目标与阶段安排}
\subsection{从关系数据库到可操作应用}
第一阶段已经能在SQL Server中保存商品，库存和订单，但从表里查到数据，还不能说明一次销售就能完整地完成。第5--9周沿着这个问题继续往下做，先检查原有关系能否表达业务，再调整需要补上的约束，最后让不同角色通过网页完成操作，并回到数据库核对结果。



\subsection{从现货销售开始}
现货销售比较适合作为这一阶段的完整流程，商品建档以后要入库，店员登记商品和数量，再确认结账，顾客最后能够查到自己的订单。这几步连在一起，价格，库存和操作身份之间的问题才会显现出来。所以应用先把这条流程做完整，原有预售和采购数据继续保留，履约和验收界面留待后续扩展。

第六周检查后保留了原有15张业务表和96个字段。网页应用需要知道当前是谁在操作，因此另外增加 \code{WebUser} 登录表，两个编号序列，五类受限数据库身份和顾客订单行级过滤，业务表本身的含义没有随界面开发改变。

\subsection{实验环境}
应用使用Python 3.12和Flask，通过pymssql连接SQL Server，运行环境由Conda管理。实验另建 \code{AnimeGoodsStoreDB_Week6App}，依次装入基础数据，迁移和应用安全结构。后面的操作主要在网页完成，SSMS用来核对同一笔业务实际保存了什么。

\section{第5周 ER 模型与业务规则检查}
\subsection{实体划分与联系表达}
这次ER建模从已经存在的15张表往回检查，重点看一张表在业务中表示什么，以及它与其他记录怎样联系。商品类别，IP和人员资料可以独立存在，一笔订单则需要通过明细才能说明买了哪些商品。按这种关系重新看表，比只按表名连线更容易发现遗漏的规则。

\Needspace{0.68\textheight}
\subsection{商品与库存的基数和参与约束}
商品有时能明确对应一个角色，有时只对应一部作品。比如群像商品不适合强行填入某个角色，所以IP和具体角色允许按实际情况填写。指定角色以后，所属IP就不能任意选择，两者的组合约束用来避免把一个角色放到另一部作品下面。每件商品仍须归属一个类别。

库存以 \code{product_id} 同时作为主码和外码，所以一件商品至多对应一条当前库存。原结构允许商品已经建档而库存行尚未建立，网页中再遇到这种情况，就要区分没有库存记录与库存为零。应用因此在新增商品的同一事务里补建零库存行，后续入库便能从明确的零开始。

\shot[0.57\textheight]{er-product.png}{商品与库存模块 ER 图}{fig:erproduct}

图中用最小和最大次数标记参与联系的范围，下划线标记主码，AK辅助标记候选码，问号和斜线分别表示可空与派生属性。这里的可售数量由实际库存减去预留量得到，查询时算出即可，另存一份反而还要在每次库存变化时同步维护。

\medskip
\Needspace{0.68\textheight}
\subsection{销售订单与人员联系}
销售订单先记录由谁收银，购买者可以是会员，也可以是散客。具体买了什么放在明细中，以 \code{order_id, product_id} 共同定位一行。同一商品再次加入同一订单时，只需修改这一行的数量，数量和登记单价也都属于这笔订单里的商品，不能只由商品编号决定。

订单与明细的关系还要结合状态理解。刚创建草稿时没有商品是正常的，到了结账时，空订单却应被拒绝。明细外码只能保证它指向的订单存在，无法独自判断订单现在能否完成。所以这条规则要继续放到结账检查中，完成后的明细也随状态变为只读。

\shot[0.57\textheight]{er-sales.png}{销售，明细与人员模块 ER 图}{fig:ersales}

收银员工存在，不等于当前登录者就有权完成收银。业务角色表保存员工职责，网页登录后还要核对员工是否启用，以及账户绑定的职责是否相符。这样订单中的经办身份才能与实际登录身份对应起来。

\medskip
\Needspace{0.72\textheight}
\subsection{预售与采购规则及原结构问题}
会员取消预订后还可能重新预订，如果把活动和会员的组合作为唯一标识，就会挡住这段需要保留的历史。因此 \code{presale_id, member_id} 不能直接作为预订表候选码。提货时还要对应原来的会员和商品，一条预订最多形成一笔提货订单，这些条件也需要连起来检查。

采购采用一场预售最多一张采购单的约定，普通补货则不关联预售活动。回看原结构时，几个问题正是出现在这种约定与实际约束之间，业务上已经说明了应当怎样做，表中却仍能保存不符合约定的记录。下面将它们与后续处理放在一起核对。

\begin{table}[H]\centering\small
\begin{tabularx}{\textwidth}{P{1.0cm}P{4.0cm}Y}
\toprule 编号 & 原结构问题 & 后续处理与保留的问题 \\
\midrule
P01 & 同一预售可能关联多张采购单。 & 第六周为非空活动编号增加筛选唯一索引，限制同活动最多一张采购单。 \\
P02 & 状态包含部分到货，与整单到货约定不一致。 & 删除 \code{PARTIAL_RECEIVED} 状态，保留明细到货数量，数量与整单状态的联动仍未强制。 \\
P03 & 完成订单可能无明细或继续被修改。 & 应用允许草稿编辑，结账拒绝空单，完成后交易内容只读。 \\
P04 & 商品建档后可能没有库存行。 & 应用新增商品时同事务创建零库存行，库存调整也处理缺失记录。 \\
P05 & 普通补货可能错误关联预售活动。 & 第六周以双向 CHECK 同时约束采购类型和活动编号。 \\
P06 & 预售会员，商品，预订状态存在跨表规则。 & 留待完整预售履约接口处理。 \\
P07 & 销售，库存与操作身份可能脱节。 & 当前现货结账通过登录身份，角色检查及统一事务实现，采购和预售链尚待扩展。 \\
\bottomrule
\end{tabularx}\caption{由 ER 反查得到的问题与落实位置}
\end{table}

这些问题有的能直接变成唯一约束或CHECK，有的却要等到一次业务操作中同时检查几张表。ER图把关系整理出来以后，接下来还要判断规则应落在哪里，才能从图上的约定走到实际能够拒绝错误数据的实现。

\clearpage
\section{第6周函数依赖与规范化取舍}
\Needspace{0.72\textheight}
\subsection{候选码与主要函数依赖}
样例中某列没有重复，还不足以说明它就是候选码。这里先根据业务约定和唯一约束确认哪些属性能标识一条记录，再检查其他属性依赖单个码，整组码，还是其中某个非主属性。这样才有依据讨论部分依赖和传递依赖。

\begin{table}[H]\centering\small
\begin{tabularx}{\textwidth}{P{4.0cm}P{3.0cm}Y}
\toprule 关系 & 主要候选码 & 依赖与分析结论 \\
\midrule
\code{ProductCategory} & 类别ID，类别名称 & 任一候选码确定描述等属性，类别名称单独存储，避免在商品行内反复更新。 \\
\code{CharacterInfo} & 角色ID，IP与角色名称组合 & 角色ID确定IP及角色名称，辅助唯一组合 \code{character_id, ip_id} 是超码，但不是最小候选码。 \\
\code{Product} & 商品ID，商品编号 & 商品码确定商品资料，在指定角色的子集内还存在角色ID决定IP的依赖，保留受控冗余。 \\
\code{Inventory} & 商品ID & 商品ID确定当前数量，预留数量，警戒值和更新时间，可售数量由查询计算。 \\
\code{SalesOrder} & 订单ID，订单编号 & 订单码确定会员，收银员及状态，非空预订编号的条件唯一性不作为全表普通候选码。 \\
\code{SalesOrderItem} & 订单ID与商品ID组合 & 整组码确定数量和成交单价，不将商品当前售价误作历史成交价的决定因素。 \\
\code{PresaleReservation} & 预订ID，预订编号 & 保留取消与重订历史，活动与会员组合不保证全表唯一。 \\
\code{PurchaseOrderItem} & 采购ID与商品ID组合 & 整组码确定采购量，成本和到货量，数量约束与状态联动需要分别判断。 \\
\bottomrule
\end{tabularx}\caption{代表性关系的候选码与依赖分析}
\end{table}

例如会员手机号有唯一约束，可以据此分析候选码，员工手机号和供应商名称则不能照搬这个结论。再看订单明细，商品名可以由商品编号确定，但买了几件，以什么单价登记，还必须知道是哪一笔订单。所以商品资料单独保存，数量和单价留在明细中。

\medskip
\Needspace{0.50\textheight}
\subsection{规范化示例与分解性质}
为了说明现有商品表与类别表为什么分开，先在理论上把它们合成一个关系。
\[
R(\mathit{product\_id},\mathit{product\_name},\mathit{category\_id},\mathit{category\_name})
\]
商品ID决定商品名与类别ID，类别ID又决定类别名。如果每件商品都带着一份类别名称，改一次类别名就要改多行。还没有商品的类别难以单独插入，最后一件商品被删除时，类别本身也可能跟着丢失，更新异常便由这种重复保存产生。

把商品资料与类别资料分开以后，得到下面两个关系。
\begin{align*}
R_1 &= (\mathit{product\_id},\mathit{product\_name},\mathit{category\_id}) \\
R_2 &= (\mathit{category\_id},\mathit{category\_name})
\end{align*}
两个关系共有类别ID，而类别ID能够决定 \(R_2\) 的全部属性，所以按类别ID连接时，可以把商品与它所属的类别重新对应起来，这个分解是无损的。原来的两组依赖分别留在 \(R_1\) 和 \(R_2\) 内部，检查它们也不需要先把表连接回来，因此依赖得到保持。

这个合并关系只用于分析。原库中的商品和类别已经分开，第六周无需再执行一次相同的拆分。规范化检查给出的结果是继续保留这部分结构，实际迁移再去处理采购约束中确实存在的问题。

\subsection{保留字段的设计理由}
\textbf{商品同时保留IP和具体角色。} 从角色能够查到所属IP，看起来商品中的IP可以省去。但群像商品只关联作品，不一定对应某个具体角色，删去IP就失去了这部分表达。另一方面，在角色非空的记录中，角色ID决定IP，角色ID却不是商品表超码，IP也不是主属性，这部分条件依赖不满足严格第三范式。这里保留两个字段，用组合外码检查它们是否匹配，同时承认需要维护这份冗余。

\textbf{保留历史值，派生值尽量查询计算。} 商品现在卖多少，与某笔订单按多少成交，是两个时点的问题。明细单价要留下本次购买的价格，目录以后调价也不能改掉这段记录。总额则由明细数量乘单价求和，可售量由实际库存减预留量得到，查询时算出就能与当前记录一致，不必再保存一份等待同步的数值。

\medskip
\Needspace{0.74\textheight}
\subsection{快照，到货数量与约束的剩余边界}
预售需求也有类似的时点区别。决定采购时参考的需求量需要留下来，之后有人取消或重订，当前需求会变，但不应覆盖当时作出决定的依据。计划采购量还可能考虑已有库存，所以它也不一定等于需求快照加额外备货量。

订金实收和退款额也按实际收付保留。如果以后明确规定实收额始终等于数量乘固定订金单价，才需要进一步讨论这部分可计算冗余。目前尚未把这种假设写成完整约束，范式分析也就保留它的适用条件。

\code{received_qty} 则需要回到当前真正实现的约束来看。若整单收货和明细状态已经严格同步，到货量有可能由状态与采购量推导出来。但现在只限制它处于零到采购量之间，尚未保证每行与整单状态联动。直接删掉这列，等于先假定了还没有完成的流程，因此这次仍保留到货量。

\begin{table}[H]\centering\small
\begin{tabularx}{\textwidth}{P{3.4cm}Y P{3.2cm}}
\toprule 选择 & 采用理由 & 还需保证什么 \\
\midrule
保留商品IP与角色 & 支持群像商品及角色商品的不同表达。 & 校验角色归属IP，避免不匹配。 \\
保留明细成交价 & 保存交易历史，允许之后调整商品售价。 & 登记时保存目录标价，结账沿用，完成后禁止修改。 \\
查询计算总额与可售量 & 数值由现有记录确定，避免重复维护。 & 查询与页面使用一致计算口径。 \\
保留需求决策快照 & 保存采购决策时依据，便于追溯。 & 区分快照与当前汇总值。 \\
暂保留采购到货量 & 保持原数据与操作含义。 & 后续实现数量与收货状态的事务联动。 \\
\bottomrule
\end{tabularx}\caption{结构保留与受控冗余的取舍}
\end{table}

所以这次保留15张业务表，把实际调整集中在采购规则上。保留一列同样需要说明理由，规范化分析既要指出依赖，也要看删改之后还能不能表达原来的业务。

\clearpage
\section{迁移方案与回归结果}
\subsection{三项结构调整及执行顺序}
先留一份未迁移的基线，再在相同数据的副本上执行调整，后面才能对照变化。脚本检查旧结构与冲突记录，满足前提后才在一个事务中建立索引和替换CHECK，冲突时列出记录并停止，保留原数据供核对。

\begin{lstlisting}[language=SQL]
CREATE UNIQUE INDEX UX_PurchaseOrder_PresaleID
ON dbo.PurchaseOrder(presale_id)
WHERE presale_id IS NOT NULL;

ALTER TABLE dbo.PurchaseOrder WITH CHECK
ADD CONSTRAINT CK_PurchaseOrder_TypePresale CHECK
((purchase_type='PRESALE' AND presale_id IS NOT NULL)
 OR (purchase_type='RESTOCK' AND presale_id IS NULL));

ALTER TABLE dbo.PurchaseOrder WITH CHECK
ADD CONSTRAINT CK_PurchaseOrder_Status CHECK
(purchase_status IN ('CREATED','ORDERED','RECEIVED','CANCELLED'));
\end{lstlisting}

普通补货没有活动编号，唯一索引因此只检查非空的 \code{presale_id}，多张补货单仍可保存。已取消的采购单也受它限制，所以同一活动最多一张采购记录，取消后另建采购单仍会被拒绝。旧约束移除与回滚由完整迁移脚本处理。

采购类型要从两个方向检查，预售采购需要活动编号，普通补货则应为空，只约束前者仍会留下错误关联。状态域同时移除部分到货，但允许哪些状态，与到货数量怎样和状态一起变化，仍是两个需要分别处理的问题。

建立约束时使用 \code{WITH CHECK} 检查既有行，脚本同时启用 \code{XACT_ABORT} 和筛选索引所需选项。旧结构不符或已迁移过时就停止，避免在错误前提下改表。

\subsection{冲突处理与实际调试}
同活动多单，类型关联不符和非法状态都会使迁移停止，原记录仍保留，已迁移的副本再次执行则返回51005。调试中还实际遇到过1934错误，原因是筛选索引要求的会话选项不一致。补齐SET并关闭 \code{NUMERIC_ROUNDABORT} 后，45项回归通过。索引建立以后，后续写入仍受这些选项影响，所以应用连接也沿用相同设置。

\medskip
\subsection{迁移前后语义保持}
\begin{table}[H]\centering\small
\begin{tabularx}{\textwidth}{P{2.7cm}P{1.2cm}Y}
\toprule 回归分类 & 通过数 & 检查内容 \\
\midrule
DATA & 15 & 全部业务表的行数及全列双向 \code{EXCEPT}，检查数据是否增减或改变。 \\
QUERY & 5 & 订单明细，金额，销量，库存等关键查询，检查业务结果是否保持。 \\
META & 7 & 旧结构与目标结构的键，索引，约束等元数据。 \\
TARGET & 1 & 新增筛选唯一索引和两项CHECK均启用且受信任。 \\
WRITE & 4 & 合法写入仍然可以执行。 \\
REJECT & 7 & 重复活动采购，类型不符，非法状态等反例被拒绝。 \\
ROLE & 6 & 既有员工角色的允许与拒绝行为。 \\
\midrule
合计 & 45 & 45项通过，0项失败。 \\
\bottomrule
\end{tabularx}\caption{第六周独立副本回归结果}
\end{table}

\begin{table}[H]\centering\small
\begin{tabularx}{\textwidth}{P{2.4cm}P{2.0cm}P{2.0cm}Y}
\toprule 核对项目 & 迁移前 & 迁移后 & 解释 \\
\midrule
SO004金额 & 124.00 & 124.00 & 两条明细汇总保持。 \\
6笔订单金额合计 & 530.00 & 530.00 & 关键查询语义保持。 \\
6件商品可售库存合计 & 42 & 42 & 数量为8，5，2，12，6，9。 \\
\bottomrule
\end{tabularx}\caption{固定样例的迁移前后对照}
\end{table}

SO004仍为124元，6笔订单合计530元，可售库存合计42件，这些结果先说明常用查询没有因迁移而改变。再结合全部表的全列双向EXCEPT，以及合法写入和非法输入的测试，才能进一步判断原数据得到保留，新规则也确实起作用。行数相同只是其中一部分依据。

\clearpage
\section{第7--9周应用设计与实现}
\subsection{从一次操作理解分层}
点下确认结账时，页面能够检查是否填了数量，却无法独自判断库存是否还够，订单是否已经完成。因此页面只把操作交给后端，业务层决定能不能做，数据访问层再执行SQL并管理连接。这样同一条规则可以在各个入口共用，不会只在某个按钮里生效。

\begin{figure}[H]\centering
\begin{tikzpicture}[font=\small,>=Stealth,
box/.style={draw=myteal,rounded corners,fill=myteal!4,text width=12.5cm,align=center,minimum height=1.2cm},link/.style={->,thick,color=myteal}]
\node[box] (a) {页面层\quad HTML、JavaScript、Flask 路由\\输入与会话、CSRF、API 调用、统一反馈};
\node[box,below=0.65cm of a] (b) {业务层\quad services.py\\身份与角色、字段规则、草稿状态、结账流程与事务范围};
\node[box,below=0.65cm of b] (c) {数据访问层\quad db.py、repository.py\\连接与释放、参数绑定、SQL、提交与回滚};
\node[box,below=0.65cm of c] (d) {数据库层\quad SQL Server\\业务表、键和约束、受限账户、顾客订单行级安全};
\draw[link] (a)--(b);\draw[link] (b)--(c);\draw[link] (c)--(d);
\end{tikzpicture}\caption{应用分层及请求流向}
\end{figure}

\begin{table}[H]\centering\small
\begin{tabularx}{\textwidth}{P{6.0cm}Y}
\toprule 模块 & 负责内容 \\
\midrule
\code{app/templates/index.html} & 页面结构，表格与对话框。 \\
\code{app/static/app.js} & 查询，表单提交，角色入口和错误展示。 \\
\code{app/__init__.py} & 路由，会话，CSRF，统一响应与异常映射。 \\
\code{app/services.py} & 权限，输入与业务状态校验，组织完整业务操作。 \\
\code{app/db.py} & 建立同事务连接，提交，回滚和关闭。 \\
\code{app/repository.py} & 参数化 SQL 与结果读取。 \\
\code{sql/01_web_security.sql} & 登录技术表，受限角色与顾客行级过滤。 \\
\bottomrule
\end{tabularx}\caption{各层在源代码中的位置}
\end{table}

\medskip
\Needspace{0.85\textheight}
\subsection{查询与写入的执行过程}
以搜索商品为例，页面将输入的名称或编号交给Flask路由，后端先确认登录身份和搜索词长度，再让Repository查询商品与库存。结果返回后才刷新表格。搜索条件经过这条路径进入数据库，而不是由页面自己组成一段SQL。

\begin{lstlisting}[language=Python]
# services.py
with self.connection(user) as repo:
    return repo.products("%" + term + "%", user["role"] in STAFF)
\end{lstlisting}
\begin{lstlisting}[language=SQL]
-- repository.py 中查询的参数绑定形式
WHERE (p.product_name LIKE %s OR p.product_code LIKE %s)
\end{lstlisting}
这里的百分号和搜索词拼成LIKE的参数值，SQL语句本身仍然固定。输入中即使带引号，也只会作为这个参数的一部分处理。商品名，数量和价格等写入值同样绑定为参数，查询与修改没有各自采用一套拼接办法。

新增商品比查询多了一层需要一起完成的修改。资料通过编号，类别，角色归属和价格检查后，商品与零库存行要在同一事务里保存，失败时一起撤销。删除则要再看库存和历史引用，已经参与交易的商品用停售保留记录，避免历史订单失去原来的商品关联。

库存调整接收商品编号，增减量和警戒值，保存前检查调整后的实际量不能低于预留量。草稿则接收会员，商品，数量和首次登记时看到的标价，结账仍要按当前身份重新核对订单状态。查询返回的是可见记录，写入返回成功以后还会重新读取结果，页面才能显示实际保存的状态。

这些操作的失败也由后端统一接收并转为页面提示。数量不足时保留草稿，店员可以调整后重试，连接失败则提示暂时不可用。响应与日志共用请求编号，后端另外保存异常细节，页面只显示当前需要采取的处理。

\medskip
\Needspace{0.72\textheight}
\subsection{身份与三层权限控制}
网页登录后先确定账户对应的员工或会员，再核对它是否启用以及承担什么职责。这个身份由会话和数据库绑定共同确定，页面不能自己选择成为店长。密码也只保存哈希，业务操作沿用已经核对过的身份。

\begin{table}[H]\centering\small
\begin{tabularx}{\textwidth}{P{2.7cm}P{2.4cm}P{2.5cm}Y}
\toprule 角色 & 商品 & 库存 & 订单 \\
\midrule
店长 \code{manager} & 查询，新增，修改，条件删除。 & 查询与调整。 & 创建草稿，编辑明细，结账，取消及草稿删除。 \\
店员 \code{clerk} & 查询。 & 在目录查看可售量。 & 创建草稿，编辑明细，结账，取消及草稿删除。 \\
库存员 \code{inventory} & 查询。 & 查询与调整。 & 不开放订单和会员查询。 \\
顾客 \code{customer} & 查询在售商品。 & 仅随目录查看可售量。 & 查询本人订单，不开放经营写入。 \\
\bottomrule
\end{tabularx}\caption{四类业务角色的主要操作范围}
\end{table}

界面初版把商品目录和库存管理都提供给了多个角色，两处却展示相近的内容，入口增加了，操作目的并没有分清。逐角色检查后，店长的商品页只维护资料与售价，库存页再处理数量，店员在目录查看可售量并进入收银，库存员则只保留库存管理。顾客不需要店铺经营统计，只浏览在售商品和本人订单。后端仍逐项检查权限，再选用 \code{auth}，\code{manager}，\code{clerk}，\code{inventory} 或 \code{customer} 受限身份连接数据库，页面上的简化没有代替这些检查。

两个顾客账号的角色相同，能查到的记录却应不同。所以顾客权限还要带上绑定的会员编号，连接中设置只读 \code{member_id} 上下文，再由SQL Server的行级安全策略过滤订单和明细。API也按会员筛选，后面的测试会分别检查页面范围和跨会员直接请求。

页面校验能及时提示输入问题，后端仍要再检查一次。价格按十进制处理并限制正值，两位小数和上限，数量则拒绝零，负数，布尔值和非整数，写请求同时校验CSRF。这样绕过表单直接提交请求时，仍然经过同样的规则。

\medskip
\Needspace{0.82\textheight}
\subsection{草稿与结账的事务边界}
订单先存为 \code{CREATED} 草稿，允许调整会员和数量，暂不扣库存。价格却不能因为改数量而跟着改变。原实现会在改数量和结账时取最新售价，但顾客是先看标价再决定购买，确认前后的金额应当一致。所以改为首次登记时核对并保存目录标价，已有明细的刷新，改数量和结账都沿用它。尚未登记时若页面价格已过期，就提示刷新核对，拒绝按另一个价格直接写入。结账成功时转为 \code{COMPLETED} 并保留记录。

结账时先锁定订单，确认它仍是可修改的草稿，会员有效且已经有明细，再按商品ID顺序检查在售状态与可售库存。关键是让所有明细共用一个事务。若每扣完一件就提交，后面的商品不足时，前面已经提交的扣减便无法随整单一起撤销。

\begin{lstlisting}[language=SQL]
UPDATE dbo.Inventory
SET on_hand_qty = on_hand_qty - %s,
    updated_at = SYSDATETIME()
WHERE product_id = %s
  AND on_hand_qty - reserved_qty >= %s;
\end{lstlisting}

库存是否足够也放进UPDATE的条件里，只有可售量满足数量时才扣减，再根据影响行数判断结果。这样读取库存以后即使还有其他订单参与竞争，写入时仍会重新受这个条件限制。商品锁按固定顺序取得，配合事务处理并发争用。

\begin{lstlisting}[language=Python]
# db.py 中事务管理的核心结构
try:
    # connection 已建立，cursor 已设置所需 SET 选项
    yield Repository(cursor)
    connection.commit()
except BaseException:
    connection.rollback()
    raise
finally:
    connection.close()
\end{lstlisting}

后面的商品不足时，前面已经执行的扣减仍留在同一事务里，因此可以一起回滚，订单继续保持草稿，登记单价也不变。这里真正需要一起完成的是整笔销售，不能只看某条UPDATE是否执行成功。业务拒绝，约束错误和执行异常最终都沿这条回滚路径返回，连接在最后关闭。

\clearpage
\section{本机业务流程与结果分析}
\subsection{商品维护与库存入库}
这次先由店长新增 \code{DEMO001}，再把名称改为初音未来徽章，售价改为120元，类别，IP和角色分别设为徽章，初音未来和初音未来。另一件 \code{DEMO002} 的售价为20元，名称实际保存成文本 \code{None}。保存后再回到目录查询，核对页面确实读到了修改后的资料，新增商品的库存则从零开始。

\begin{figure}[H]\centering
\begin{minipage}{.47\linewidth}\centering\includegraphics[width=\linewidth,height=.29\textheight,keepaspectratio]{product-create.png}\end{minipage}\hfill
\begin{minipage}{.47\linewidth}\centering\includegraphics[width=\linewidth,height=.29\textheight,keepaspectratio]{product-edit.png}\end{minipage}
\caption{商品新增与修改表单，局部截图}
\end{figure}

接着换库存管理员入库，两件商品分别增加10件和5件，警戒值设为3和2，查询得到可售量合计15。商品建档与入库在这里是两步，库存员可以处理数量，却不能顺便修改商品售价或查看订单。店员需要的只是目录中的售价和可售量。

\shot[0.22\textheight]{stock-before.png}{入库后两件商品库存为10和5，局部截图}{fig:stockbefore}

\clearpage
\subsection{草稿编辑，结账与数据库核对}
店员为会员 \code{M001} 建立 \code{WEB0000010015} 草稿，加入两件 \code{DEMO001} 和一件 \code{DEMO002}，金额为260元。先把第一行改为1件并保存，金额随之变成140元，再恢复为2件。这一步同时检查数量修改与重新汇总，草稿期间库存仍保持原数，最后才确认结账。


确认以后，订单变为已完成，明细不能再编辑。两行金额分别为240元和20元，总额260元，库存则从10和5变成8和4，正好分别减少2件和1件。页面上的金额与数量变化能够对应起来。

\shot[0.16\textheight]{completed.png}{结账后保存成交价格与订单明细，局部截图}{fig:completed}

页面显示成功之后，再到SSMS查询同一笔订单。关联商品，库存和明细后，状态为 \code{COMPLETED}，会员ID为1，数量与单价也都和网页一致，说明这些变化已经写进数据库，而不只是页面上的提示。

\shot[0.19\textheight]{sql-check.png}{数据库中的库存与成交明细，局部截图}{fig:sql}

\clearpage
\subsection{两位顾客的本人订单查询}
登录 \code{customer1}，列表中能看到刚完成的260元订单，以及SO004和SO003。它可以打开明细核对，却不能修改商品，库存或订单，收银完成后留下的记录以只读形式提供给顾客。

\shot[0.25\textheight]{customer1.png}{顾客1可以查看会员1的订单，局部截图}{fig:customer1}

换成 \code{customer2} 后，列表只有会员2的SO002，金额90元，刚才那笔260元订单不再出现。两个账号使用同一套顾客界面，结果却随绑定会员改变，这里区分的是本人记录的范围，而不只是有没有进入订单页的权限。跨会员直接访问还会在后面的集成测试中检查。

\shot[0.18\textheight]{customer2.png}{顾客2只显示会员2的订单，局部截图}{fig:customer2}

\clearpage
\subsection{库存不足与整单回滚}
接着另建 \code{WEB0000010016}，第一行买1件 \code{DEMO001}，第二行买10件 \code{DEMO002}，金额320元。库存分别为8和4，所以第一行足够，第二行不足。特意让不足出现在后面，才能检查前面已经处理的商品是否也会撤销。确认后返回 \code{INSUFFICIENT_STOCK} 和409，页面提示库存不足，订单仍是草稿。

\shot[0.31\textheight]{checkout-error.png}{库存不足时拒绝结账并保留草稿，局部截图}{fig:failure}

再查库存，结果仍是8和4。结账按商品ID顺序处理，第一件本来可以扣减，但第二件失败后，它的扣减也被撤销了。这里比单独验证一件商品不足更能说明事务范围，失败不会留下部分商品已扣减而订单尚未完成的状态。日志也能按页面上的请求编号找到这次操作。

\shot[0.14\textheight]{rollback-stock.png}{失败后两件商品库存保持8和4，局部截图}{fig:rollback}

这笔失败草稿随后确认删除，再查询时该订单编号已不存在。草稿还没有形成成交记录，可以修改后重试或删除，完成订单则保留原交易内容。其他删除条件，引用冲突和越权情况由自动化测试继续检查。

\clearpage
\section{集成测试与结果分析}
最终执行pytest得到48项通过，0失败，0错误，0跳过。其中42项连接真实SQL Server检查应用行为，另外6项检查密码重设辅助脚本。

\begin{table}[H]\centering\small
\begin{tabularx}{\textwidth}{P{3.1cm}Y P{1.2cm}}
\toprule 测试范围 & 代表性检查 & 结果 \\
\midrule
登录与请求保护 & 登录退出，未登录访问，密码哈希与后端身份，CSRF与非法JSON。 & 通过 \\
商品与库存 & 商品CRUD，自动零库存，重复约束回滚，非法价格，预留保护，有库存和有引用商品删除拒绝。 & 通过 \\
订单与事务 & 草稿及明细CRUD，空单拒绝，目录标价核对，草稿锁价，多商品结账，历史价保持，库存不足整单回滚。 & 通过 \\
角色与顾客隔离 & API越权拒绝，顾客跨会员访问与SQL RLS，受限SQL账户直接执行拒绝。 & 通过 \\
并发与故障 & 并发结账不超卖，数据库连接失败与日志，页面安全响应头。 & 通过 \\
密码重设辅助 & 生成哈希匹配且不含明文，非法输入拒绝，已有输出文件保留。 & 通过 \\
\bottomrule
\end{tabularx}\caption{本机48项测试的范围与结果}
\end{table}

价格规则另外做了两种检查。商品按25元登记以后把目录调为30元，刷新和改数量仍保留25元，两件结账50元，库存由10减到8。还没登记的商品若使用过期页面标价，则由集成测试确认拒绝写入，提示重新核对。前者保留已登记价格，后者避免新登记时无提示换价，处理时点并不相同。

参数化查询直接使用带引号和SQL片段的搜索词，返回空列表以后，商品仍可正常查询。这个输入没有作为SQL执行，再结合Repository的参数绑定方式，检查查询条件与语句结构确实分开。

权限测试则绕过隐藏按钮，直接让顾客新增商品，店员调整库存，库存员查询订单，均得到拒绝。再用受限SQL账户执行越权语句，返回229错误。顾客跨会员访问同时检查API和数据库过滤，所以页面列表正确只是其中一项观察。

断连测试把连接指向不可用端口，页面得到503，并能按请求编号找到日志。Windows复跑时曾在读取这份日志处失败，实际是UTF-8日志被按默认编码读取，明确读取编码后全量复跑48项通过。这里需要区分连接失败已经正确处理，与测试读取记录本身出了问题。

并发用例设置实际库存3件，预留1件，让两个订单各买2件，结果只有一笔成功，另一笔返回库存不足，最终实际库存与预留量均为1。这个结果验证了所设置的争用场景，没有把两次购买都算作成功，尚不能代替更大规模的压力测试。

\clearpage
\section{分工，AI辅助与实验总结}
\subsection{小组分工与辅助过程}
郑立峰负责第5--6周的ER模型，规则检查，依赖与范式分析，以及采购约束迁移和回归材料。李维政在这个结构上完成第7--9周应用，连接数据库，区分角色操作，处理结账事务与异常，并进行集成测试和本机流程截图，再将各周结果整理进报告。最终提交前仍需共同核对图，代码和实际结果是否一致。

AI参与了规则分析，代码和测试组织，连接排查，以及报告的整理。业务范围和字段取舍由小组确认，本机截图记录实际的角色切换，数量修改，结账和失败操作。自动测试的运行结果另行保留，报告里的价格，金额与库存都对应已有截图或查询结果，未用预设演示值替换实际观察，过程见 \code{ai_log.md}。

\subsection{实验认识与当前不足}
回看这一阶段，表结构和完整流程之间还有不少需要补上的判断。外码能保证明细找到订单，却不能判断空订单能否结账，角色存在也不等于当前账户就有收银权限。把这些关系真正放进一次操作以后，才能分清哪些交给数据库约束，哪些需要由业务层继续处理。

价格规则的修改也让历史值的含义更具体了。目录价说明商品现在怎样定价，登记单价说明这笔购买接受了什么价格，两者在不同时间发生变化并不矛盾。已有明细需要保留约定，新登记又要核对页面所见价格，单纯在结账时取一次最新售价，反而会破坏确认前后的对应关系。

库存失败案例则给出了更直接的核对方式，第一件够卖，第二件不足，最后两件库存都没有变化，才能说明整笔销售确实回到了原来的状态。一次成功扣减本身看不出这个区别，所以测试中仍要保留会失败的输入，并核对失败之后留下了什么。

现货流程目前已有对应的页面，数据库结果和测试依据。预售履约与采购验收还没有沿同样方式完成，到货数量和整单状态的联动也仍需继续实现，这部分不能从当前现货测试的通过结果直接推出。

\end{document}
